
	

%%%%%%%%%%%%%%%%%%%%%%%%%%%%%%%%%%%%%%%%%%%%%%%%%%%%%%%%%%%%%%%%%%%%%%%%%%%%%%%
%                                                                             %
%  TITEL:  Abschlussarbeit 													%		%	                    														%
%          --- Hauptdokument ---                                              %
%                                                                             %
%  DATEI:  AbschlussLatex.tex                                                 %
%                                                                             %
%  DATUM:  01.09.2020                                                         %
%                                                                             %
%%%%%%%%%%%%%%%%%%%%%%%%%%%%%%%%%%%%%%%%%%%%%%%%%%%%%%%%%%%%%%%%%%%%%%%%%%%%%%%


%--- Vorspann ----------------------------------------------------------------%

\documentclass[fontsize = 12pt,					%Schriftgr��e
               paper = a4,						%Papierformat
               headings = small,					%Gr��e der �berschriften
               open=right,						%Abschnitte beginnen rechts
               cleardoublepage = empty,			%leere Seiten ohne Kopfzeile
			   BCOR = 10mm,						%Binde Korrektur
               captions = tableheading,			%Tabellen mit �berschriften
               bibliography = totoc,				%Literatur- ins Inhaltsverzeichnis
               listof = totoc,					%Verzeichnisse ins Inhaltsverzeichnis (totoc)
               twoside = false,					%doppelseitiges Layout
               american]							%Dokumentklasse
              {scrbook}							

 
%--- Kopf- und Fu�zeilen -----------------------------------------------------%

\usepackage{scrpage2}
\usepackage{scrhack}
\pagestyle{scrheadings}
\automark[section]{chapter}
\addtokomafont{pageheadfoot}{\normalfont\sffamily}
\addtokomafont{pagenumber}{\normalfont\sffamily\bfseries}
\setheadsepline{0.4pt}
\clearscrheadfoot
\cehead[]{\leftmark}
\cohead[]{\rightmark}
\ohead[]{\pagemark}
\ofoot[\pagemark]{}


%--- Tabellen- und Abbildungsbeschreibungen ----------------------------------%

\addtokomafont{caption}{\small}
\addtokomafont{captionlabel}{\sffamily\bfseries}


%--- Gliederungstiefe (Bezifferung/Inhaltsverzeichnis) -----------------------%

\setcounter{secnumdepth}{2}
\setcounter{tocdepth}{2}


%--- Eingabekodierung --------------------------------------------------------%

%------ Linux ----------------------------------------------------------------%
%\usepackage[latin1]{inputenc}

%------ Windows --------------------------------------------------------------%
\usepackage[ansinew]{inputenc}


%--- Schrift -----------------------------------------------------------------%
\usepackage{lmodern}
\usepackage[T1]{fontenc}
\usepackage{calligra}
\usepackage{wedn}
%\usepackage{fontspec}
%\setmainfont{QTHandwriting}
%\setmainfont{QTFlorencia}

% Spezielles Seitenlayout:
\usepackage[Sonny]{fncychap}		%Sonny/Lenny/Glenn/Conny/Rejne/Bjarne
\ChTitleVar{\LARGE\bfseries}
\ChNameVar{\large\bfseries}

%--- Math stuff --------------------------------------------------------------%
\usepackage{amsmath}
\usepackage{amssymb}
\usepackage{pifont}
\newcommand{\greencheckmark}{{\color{Green}\checkmark}}
\newcommand{\redxmark}{{\color{red}\ding{53}}}
\newcommand{\xmark}{{\color{red}\ding{53}}}

\usepackage{bm}
\usepackage{mathtools}
\newcommand{\myM}[1]{\bm{\mathit{#1}}} %Bold, italic, capital expression of matrices; vectors bold, italic, small
\DeclareMathOperator*{\argmin}{argmin}
\newcommand{\bfun}{\myM{f}}
\newcommand{\bg}{\myM{g}}
\newcommand{\bk}{\myM{k}}
\newcommand{\bl}{\myM{l}}
\newcommand{\bn}{\myM{n}}
\newcommand{\bp}{\myM{p}}
\newcommand{\bq}{\myM{q}}
\newcommand{\btau}{\myM{\tau}}
\newcommand{\bt}{\myM{t}}
\newcommand{\bu}{\myM{u}}
\newcommand{\bv}{\myM{v}}
\newcommand{\bw}{\myM{w}}
\newcommand{\bx}{\myM{x}}
\newcommand{\bK}{\myM{K}}
\newcommand{\bQ}{\myM{Q}}
\newcommand{\bV}{\myM{V}}
\newcommand{\bdq}{\dot{\myM{q}}}	
\newcommand{\bddq}{\ddot{\myM{q}}}
\newcommand{\du}{\delta\myM{u}}
\newcommand{\dx}{\delta\myM{x}}



%--- Weitere Pakete ----------------------------------------------------------%
\usepackage{setspace}            % Zeilenabstand einstellbar
\usepackage[usenames, dvipsnames,table,xcdraw]{xcolor} %tabellen mit farben
\definecolor{pptblue1}{rgb}{0.858,0.921,0.972}
\definecolor{pptblue2}{rgb}{0.709,0.8,0.894}
\definecolor{pptblue3}{rgb}{0.427,0.576,0.741}
\definecolor{pptblue4}{rgb}{0.211,0.407,0.627}
%\usepackage[hidelinks,colorlinks=true,linkcolor=cyan,citecolor=cyan]{hyperref}
%\usepackage[colorlinks, linkcolor = pptblue4, citecolor = pptblue4, filecolor = black, urlcolor = pptblue4]{hyperref} %Add linktocpage to have only sites clickable in ToC
\usepackage[colorlinks, linkcolor = black, citecolor = black, filecolor = black, urlcolor = pptblue4]{hyperref}

\usepackage{subcaption}
\usepackage{graphicx}
\usepackage{epstopdf}
\usepackage[export]{adjustbox} %f�r die gr��e von bildern
\usepackage{listings} %statt verbatim?
\usepackage{bibgerm}
\usepackage{float}        % u.a. genaue Plazierung von Gleitobjekten mit H
\usepackage{textcomp} 		%f�r \textregistred in �berschrift
\usepackage{tabularx}
\renewcommand{\bf}{\normalfont \bfseries} 
\usepackage{fancyvrb}  %zur Darstellung von Quelltexten
\usepackage[main=american,ngerman]{babel}
\usepackage[babel]{csquotes}
\usepackage{siunitx}  % degree symbol by \ang{degrees}
\usepackage[acronym,nonumberlist,nopostdot,automake,style=super,nogroupskip]{glossaries} % acronym option f�r seperates Glossar und Abk�rzungsliste -> sonst zusammen
\clubpenalty = 10000             % Schusterjungen und Hurenkinder vermeiden
\widowpenalty = 10000            % siehe dazu 
\displaywidowpenalty = 10000     % http://www.cs.uu.nl/wais/html/na-dir/de-tex-faq/part5.html
\usepackage[all]{nowidow}
\usepackage{listings}
%\lstset{frame=tb,
%  language=C++,
%  breaklines=true,
%  showstringspaces=false,
%  columns=flexible,
%  numbers=none,
%  tabsize=3
%}
\definecolor{dkgreen}{rgb}{0,0.6,0}
\definecolor{dred}{rgb}{0.545,0,0}
\definecolor{dblue}{rgb}{0,0,0.545}
\definecolor{lgrey}{rgb}{0.9,0.9,0.9}
\definecolor{gray}{rgb}{0.4,0.4,0.4}
\definecolor{darkblue}{rgb}{0.0,0.0,0.6}
\lstset{language=C++,
                basicstyle=\ttfamily,
                breaklines=true,
                tabsize=3,
                keywordstyle=\color{dblue}\ttfamily,
                stringstyle=\color{orange}\ttfamily,
                commentstyle=\color{dkgreen}\ttfamily,
                morecomment=[l][\color{dred}]{\#}
}


%Literatur
\usepackage{cite}                % Sortierte und zusammengefasste Zitatnummern 
\usepackage[numbers]{natbib}     % Refernz mit nummern UND \citauthor{} m�glich 
\typearea[current]{current}      % Neuberechnung des Satzspiegels mit alten Werten nach ?nderung von Zeilenabstand,etc
\renewcommand{\arraystretch}{1.2} %Tabellenh�he vergr��ern
\usepackage{subcaption}			%For subtables (table a and table b automatically)

%\renewcommand{\thesubfigure}{\arabic{subfigure}} %makes subcaptions in figs/tbls 1,2,3 instead a,b,c
\renewcommand{\glsnamefont}[1]{\textbf{#1}} %glossary entries bold
\setlength{\glsdescwidth}{5in} %gloassary terms without line breaks
\makeglossaries					%Erstelle generell Glossare

%%% abbreviations:
\newacronym[longplural={degrees of freedom}]{DOF}{DOF}{degree of freedom}
\newacronym{COG}{COG}{center of gravity}

%%% nomenclature:
%\newglossaryentry{space_joint}{name={$\myM{q}=(q_1,q_2)^T$},description={minimal coordinates in joint space}}



		%include .tex data with definitions 
\usepackage[capitalise]{cleveref}
\captionsetup[subfigure]{subrefformat=simple,labelformat=simple}
%%%%%%%% Hints for glossary:
% call in text with \gls{label} or \glspl for plural (Gls/Glspl for capital)


%--- Hauptdokument -----------------------------------------------------------%

\begin{document}
\renewcommand*{\chapterpagestyle}{empty}

\frontmatter

%-----Titelblatt, Abstract, Kurzfassung, Inhaltsverzeichnis, Figuren, Tabellen,Terms (r�m. Ziffern)-----
%-----------------------------------------------------------------------------%
%                                                                             %
%    T I T E L S E I T E                                                      %
%                                                                             %
%-----------------------------------------------------------------------------%
%
\thispagestyle{empty}
\hypersetup{pageanchor=false}
%
\begin{flushleft}
University of Duisburg-Essen\\[1ex]
Faculty of Engineering\\[1ex]
Chair of Mechatronics \\[2ex]
\end{flushleft}
%
\vfill\vfill\vfill\vfill
%
\begin{center}
\Large
\textsf{\textbf{Highly-Dynamic Movements of a Humanoid Robot Using Whole-Body Trajectory Optimization}}\\[5ex]
\large
\textbf{Master Thesis}\\[1ex]
\textbf{Maschinenbau (M.Sc.)}\\[5ex]
Julian E�er\\[1ex]
Student ID: 3015459\\[5ex]
\begin{tabular}{ll}
\textbf{First examiner} & Prof. Dr. Dr. h.c. Frank Kirchner (DFKI)\\[1ex]
\textbf{Second examiner} & Dr.-Ing. Tobias Bruckmann (UDE)\\[1ex]
\textbf{Supervisor} & Dr. rer. nat. Shivesh Kumar (DFKI)\\[1ex]
\textbf{Supervisor} & Dr. Olivier Stasse (LAAS-CNRS)\\[1ex]
%\textbf{Supervisor} & Dr. Carlos Mastalli (University of Edinburgh)\\[1ex]
\end{tabular}\\[7ex]
\normalsize
September 22, 2020
\end{center}
%
\vfill\vfill\vfill\vfill\vfill\vfill
%
\begin{center}
	\includegraphics[width=5.2cm]{img/logo_ude}\hspace{20pt}
	\includegraphics[width=8cm]{img/logo_dfki.png}
\end{center}
\addchap*{Declaration}


\thispagestyle{empty}

This study was carried out at the Robotics Innovation Center of the German Research Center for Artificial Intelligence (DFKI) in the Advanced AI Team on Mechanics \& Control under the direction of Dr. Shivesh Kumar. 
\vspace*{5.2cm}\\
I declare that this thesis was composed by myself, that the work contained herein is my own except where explicitly stated otherwise  in the text, and that this work has not been submitted for any other degree or professional qualification except as specified. 

\bigskip
\bigskip 
\bigskip 
\bigskip
\bigskip

\begin{flushright}
\begin{tabular}{c}
Bremen, September 18, 2020\\
\bigskip \\\\
\hrulefill\hrulefill\hrulefill\hrulefill\hrulefill\hrulefill \\
Julian E�er\\
\hspace{6cm} \\
\end{tabular}
\end{flushright}



\thispagestyle{empty}
\chapter*{Abstract}

In order to further close the gap between robots and their natural counterparts, current research is driving towards exploiting the natural dynamics of robots. This requires a rethinking about the way we control robots for moving in a more dynamic, efficient and natural way. 

Dynamic bipedal locomotion is a challenging problem and remains an open area of research. A key characteristic is the decoupling between the center of mass and the multi-body dynamics. Particular difficulties arise from effective underactuation, the mechanism complexity, as well as nonlinear and hybrid dynamics. 

A common approach is to decompose this problem into smaller sub-problems that are solved sequentially. Many state-of-the-art frameworks rely on trajectory optimization based on reduced centroidal dynamics, which is transferred to a whole-body trajectory via feedback linearization using Inverse Kinematics (IK) or Inverse Dynamics (ID). Recent research indicates that solving this mapping with a local optimal control solver, namely Differential Dynamic Programming (DDP), instead of IK/ID, produces more efficient motions, with lower forces and impacts. 

This master thesis contributes to the research field of dynamic bipedal locomotion by applying, evaluating and extending DDP-based whole-body trajectory optimization, pursuing three objectives: 
First, we develop an approach to constrain the optimal control problem allowing DDP to produce inherently balanced motions. 
Second, we evaluate our approach for quasi-static and dynamic motions in a real-time physics simulation and in real-world experiments on the lightweight and biologically inspired RH5 humanoid robot.
Third, we examine the limits of the derived whole-body planning approach and the system design via solving highly-dynamic movements. 

\vfill\vfill
\textbf{Keywords:} Humanoid Robots, Dynamic Bipedal Walking, Motion Planning, Multi-Contact Optimal Control, Differential Dynamic Programming, Whole-Body Trajectory Optimization  
 












\thispagestyle{empty}
\chapter*{Kurzfassung}

Die Bewegungsplanung f{\"u}r Laufroboter ist eine anspruchsvolle Aufgabe und bleibt ein aktives Forschungsgebiet. Besondere Schwierigkeiten ergeben sich aus der effektiven Unteraktuierung, der Komplexit{\"a}t des Mechanismus sowie der nichtlinearen und hybriden Dynamik. Ein {\"u}blicher Ansatz besteht darin, das Problem in kleinere Teilprobleme zu zerlegen, die nacheinander gel{\"o}st werden. Neue Forschungsergebnisse zeigen, dass Algorithmen aus dem Bereich der optimalen Regelung, insbesondere Differenzielle Dynamische Programmierung (DDP), effizientere Bewegungen mit geringeren Kr{\"a}ften und St{\"o}{\ss}en erzeugen.

Diese Masterarbeit leistet einen Beitrag in diese Richtung, indem eine DDP-basierte Ganzk{\"o}rper-Trajektorienoptimierung angewendet, bewertet und erweitert wird, wobei drei Ziele verfolgt werden. Erstens entwickeln wir eine Methode zur Beschr{\"a}nkung von DDP-artigen Algorithmen, um stabile Bewegungspl{\"a}ne zu erzeugen. Zweitens wird der vorgeschlagene Ansatz zur Bewegungsplanung f{\"u}r quasi-statische und dynamische Bewegungen in einer physikalischen Echtzeitsimulation und in realen Experimenten mit dem leichten und biologisch inspirierten humanoiden Roboter RH5 evaluiert. Schlie{\ss}lich werden die Grenzen des Ansatzes und des Systemdesigns anhand hochdynamischer Bewegungen untersucht.

\vfill
\noindent\textbf{Stichworte:} Bewegungsplanung, Differenzielle Dynamische Programmierung , Dynamisches Bipedales Laufen, Ganzk{\"o}rper-Trajektorienoptimierung, Humanoide Roboter, Optimale Regelung

\thispagestyle{empty}
\chapter*{Acknowledgments}
 
%Supervisors
I would like to thank my first supervisor and institute director Prof. Frank Kirchner for giving me the opportunity to work on this interesting topic at the DFKI Robotics Innovation Center. Then, I would like to thank my second supervisor Dr. Tobias Bruckmann (Research Group Leader Robotics, Chair of Mechatronics, University of Duisburg-Essen) for providing his continuous guidance and constructive feedback during this thesis.   

%Mentors
I would like to express my deepest gratitude to my main mentor and team leader Dr. Shivesh Kumar for his guidance, the numerous inspiring discussions we had throughout the course of this work, and his infectious enthusiasm for everything related to dynamics and many more surrounding topics. Thanks also goes to Dr. Carlos Mastalli (Postdoctoral researcher, University of Edinburgh, United Kingdom) and Dr. Olivier Stasse (Postdoctoral researcher, LAAS-CNRS, France) who supported me with many ideas and critical but constructive questions related to numerical optimization and humanoid robotics, respectively. 
%Special thanks to my work colleagues Vinzenz Bargsten, Heiner Peters and Daniel Harnack for the fruitful collaboration.   

%Friends, Family, monique
Finally, I want to also thank all my friends I met along the way, in the Lower Rhine area, Duisburg/Essen, Dortmund, %M{\"u}nchen, 
Oberpfaffenhofen, Bremen and elsewhere. You have each made this journey an unforgettable experience. Last but not least, I want to thank my family (especially my mother) and my girlfriend Monique (one of the bravest person I ever met) for your supportive stronghold and I am forever grateful to have you in my life. 


\bigskip

\begin{flushright}
{\Large Julian} \\
%Julian \\
\bigskip
September 18, 2020
\end{flushright}










   
\tableofcontents    
\printglossary[title=List of Symbols]
\printglossary[type=\acronymtype,title=Acronyms]	         


%-----Hauptkapitel-----
\mainmatter             % numerische Ziffern
\clearpage
\pagestyle{scrheadings}  % Kopfzeilen
\thispagestyle{scrheadings}  % Kopfzeilen

%-----------------------------------------------------------------------------%
%                                                                             %
%    K A P I T E L   1                                                        %
%                                                                             %
%-----------------------------------------------------------------------------%

\chapter{Introduction}\label{c1}

\section{Motivation}
\section{Related Work}
%\subsection{Model Induced Design} (Passive Dynamic Walking, Template Models)
%\subsection{Traditional Legged Locomotion Planning}
%\subsection{Trajectory Optimization} (Sequential, Simultaneous OR Direct, Indirect, DDP)
%\subsection{Series-Parallel Hybrid Robots}


\section{Contribution and Structure}


\section{RH5 Humanoid Robot}
The derived approaches for constrained \gls{DDP} have been tested both in simulation and real-world experiments on a full-size humanoid robot. RH5 is a lightweight and biologically inspired humanoid that has recently been developed at DFKI Robotics Innovation Center\cite{peters2017konstruktion}.

The RH5 humanoid robot (see \cref{img:rh5_robot}) is designed to mimic the human anatomy with a total size of 200cm, a weight of 62kg and a total of 32 \gls{DoF}. The two legs account for 12 \gls{DoF}, the torso and neck kinematics each for three and the arms and grippers of the robot for 16 \gls{DoF}. In order to achieve a high dynamic performance, the robot's design follows a series-parallel hybrid approach. Consequently,linkages and parallel mechanisms are utilized in most of the robots joints, e.g. the hip-flexion-extension, knee, ankle, torso and wrist. A comparison of RH5 with other state of the art humanoid robots revealed several advantages of this design approach, including better maximum velocity and torque of the ankle as well as an advantageous weight of the lower leg \cite{kumar2020survey}. The interested reader can find a comprehensive introduction on series-parallel hybrid robots in \cite[Ch.2]{kumar2019modular}. 

\begin{figure}[h!]
\centering	
\includegraphics[width=.3\textwidth]{img/rh5_robot.png}
\caption{The recently presented RH5 is a lightweight and biologically inspired humanoid used as experimental platform within this thesis.}
\label{img:rh5_robot}
\end{figure} 

% Add to Appendix: 
% - Gelenkwinkel
% - Gelenkachsen Darstellung
% - URDF (abstract and full)

%-----------------------------------------------------------------------------%
%                                                                             %
%    K A P I T E L   2                                                        %
%                                                                             %
%-----------------------------------------------------------------------------%

\chapter{Background: Dynamics of Legged Robots}\label{c2}





%-----------------------------------------------------------------------------%
%                                                                             %
%    K A P I T E L   3                                                      %
%                                                                             %
%-----------------------------------------------------------------------------%

\chapter{Contact Stability Constrained DDP}\label{c3}
\section{Idea: Constraining the CoP of each Contact Surface}

\section{Center of Pressure Constraints}
\subsection{Center of Pressure (CoP) Conditions}
Conditions for \gls{CoP} inside convex hull:
\begin{align}
\begin{split}
Pitch:& -X \leq C_x \leq X \\
Roll:& -Y \leq C_y \leq Y
\end{split}
\end{align}
Represented by the four inequality conditions:
\begin{align}
\begin{split}
-X &\leq C_x \\
C_x &\leq X \\
-Y &\leq C_y \\
C_y &\leq Y \\
\end{split}
\end{align}

\subsection{CoP Computation}
\begin{equation}
\myM{p}_{CoP} = \dfrac{\myM{n}\times\myM{M}}{\myM{F}\cdot\myM{n}} = \dfrac{\begin{bmatrix} n_x \\ n_y \\ n_z \end{bmatrix} \times \begin{bmatrix} t_x \\ t_y \\ t_z \end{bmatrix}}{\begin{bmatrix} f_x \\ f_y \\ f_z \end{bmatrix} \cdot \begin{bmatrix} n_x \\ n_y \\ n_z \end{bmatrix}} = 
\begin{bmatrix} n_yt_z - n_zt_y \\ n_zt_x-n_xt_z \\ n_xt_y-n_yt_x \end{bmatrix}\cdot \dfrac{1}{f_xn_x+f_yn_y+f_zn_y}
\end{equation}

\subsection{CoP Constraints: General Case}
These conditions are be implemented via: 
\begin{equation}
\begin{bmatrix} -Yn_0 & -Yn_1 & -Yn_2 & -n_2 & 0 & n_0 \\
-Yn_0 & -Yn_1 & -Yn_2 & n_2 & 0 & -n_0 \\
-Xn_0 & -Xn_1 & -Xn_2 & 0 & n_2 & -n_1 \\
-Xn_0 & -Xn_1 & -Xn_2 & 0 & -n_2 & n_1 \\ \end{bmatrix} \cdot
\begin{bmatrix} f^x \\ f^y \\ f^z \\ \tau^x \\ \tau^y \\ \tau^z \end{bmatrix} \leq
\begin{bmatrix} 0 \\ 0 \\ 0 \\ 0 \end{bmatrix}
\end{equation}

\subsection{CoP Constraints: Special Case of Horizontal Floor}
For the case of horizontal floor, with the according normal vector $\myM{n}=[0,0,1]$, the \gls{CoP} can be computed to: 
\begin{equation}
\myM{p}_{CoP} = \begin{bmatrix} -t_y/f_z \\ t_x/f_z \\ 0 \end{bmatrix}
\end{equation}
Represented by the four inequality conditions:
\begin{align}
\begin{split}
\tau^x &\leq Yf^z \\
-\tau^x &\leq Yf^z \\
\tau^y &\leq Yf^z \\
-\tau^y &\leq Yf^z 
\end{split}
\end{align}
Briefly summarized as:
\begin{align}
\begin{split}
\mid\mid\tau^x\mid\mid &\leq Yf^z \\
\mid\mid\tau^y\mid\mid &\leq Xf^z 
\end{split}
\end{align}
These conditions are be implemented via: 
\begin{equation}
\begin{bmatrix} 0 & 0 & -Y & 1 & 0 & 0 \\
0 & 0 & -Y & -1 & 0 & 0 \\
0 & 0 & -X & 0 & 1 & 0 \\
0 & 0 & -X & 0 & -1 & 0 \end{bmatrix} \cdot
\begin{bmatrix} f^x \\ f^y \\ f^z \\ \tau^x \\ \tau^y \\ \tau^z \end{bmatrix} \leq
\begin{bmatrix} 0 \\ 0 \\ 0 \\ 0 \end{bmatrix}
\end{equation}





\subsection{Backup: Friction Cone constraints}
\begin{align}
\begin{split}
\mid\mid f^x\mid\mid &\leq \mu f^z \\
\mid\mid f^y\mid\mid &\leq \mu f^z \\
f^z &> 0
\end{split}
\end{align}
For the case of four edges of the linear approximation of the friction cone, the equations become:
\begin{equation}
\begin{bmatrix} 1 & 0 & -\mu \\
-1 & 0 & -\mu \\
0 & 1 & -\mu \\
0 & -1 & -\mu \\
0 & 0 & -\mu \\ \end{bmatrix} \cdot
\begin{bmatrix} f^x \\ f^y \\ f^z \end{bmatrix} \leq
\begin{bmatrix} 0 \\ 0 \\ 0 \\ 0 \\ 0 \end{bmatrix}
\end{equation}

\section{Integration into the Crocoddyl Framework}
\subsection{Inequality Constraints by Penalization}
%CostModelSum, Residual, Activation
\subsection{Computing the Residual}
\subsection{Unittesting Numerical Differentiation} 





%-----------------------------------------------------------------------------%
%                                                                             %
%    K A P I T E L   4                                                        %
%                                                                             %
%-----------------------------------------------------------------------------%

\chapter{Experimental Validation}\label{c4}




%-----------------------------------------------------------------------------%
%                                                                             %
%    K A P I T E L   5                                                        %
%                                                                             %
%-----------------------------------------------------------------------------%

\chapter{Highly-Dynamic Movements}\label{c5}

\section{Formulation of the Optimization Problem}
\section{Simulation Results for Increasing Task Complexity}
\subsection{A Simple Vertical Jump}
\subsection{Forward Jumping Over Multiple Obstacles}
\section{Identification of Limits in System Design}
\subsection{Analysis of Joint Limits}
\subsection{Analysis of Torque Limits}



%-----------------------------------------------------------------------------%
%                                                                             %
%    K A P I T E L   6                                                        %
%                                                                             %
%-----------------------------------------------------------------------------%

\chapter{Validation of Planned Motions}\label{c6}
This chapter presents a two-step validation of the physical consistency of the planned motions where an online stabilization is applied to track the optimal trajectories. At first, we prove the stability of the motions in a real-time physics simulation. Following this, we explore the feasibility of the motions in real-world experiments on a full-size humanoid robot.    

\section{Validation in Real-Time Physics Simulation}\label{sec:OnlineSimulation}
This section investigates the stability of the planned motions in a physics simulator by applying an online stabilization based on joint space position control to track the trajectories obtained from \gls{OC}. To begin with, an overview of the simulation setup is given, then the tracking results of the planned motions are discussed. 

\subsection{Simulation Setup}
The optimal motions are tested in the dynamical simulation environment PyBullet \cite{coumans2016pybullet}. PyBullet is an open-source framework for robotics simulation that allows fast computation of rigid-body dynamics along with collision detection. 
%The focus of this investigation is to minimize the simulation-to-reality ('sim-to-real') gap, which is, despite constantly improving robot models, still a major problem in real-world experiments. 
The motion tracking in the simulator is set up in a similar way the motions would be tested on a real robot, namely interpolating the trajectories and closing the loop on joint position level.

The control loop on the real system (see \cref{sec:OnlineExperiments}) is running at a frequency of 1 kHz. The generated bipedal walking variants and highly-dynamic movements presented in \crefrange{c4}{c5} are generated with an discrete \gls{OC} formulation at 30 Hz and 10 Hz, respectively. To this end, the optimal trajectories are interpolated with a cubic spline in order to realize an up-scaling of the reference data to 1 kHz.  

The control architecture consists of a simple PD-controller on joint space level. The controller is supposed to track both position and velocity reference trajectories with the standard PyBullet parameters in a real-time loop running at 1 kHz. In this real-time simulation, the same URDF robot model is used as in the Crocoddyl framework, including the same maximal motor torques. Furthermore, the parameters of the rigid contact models have been aligned between both frameworks to ensure comparability.

\subsection{Results}
Following up, we investigate the capabilities of the presented control architecture for tracking the planned motions. To this end, we study the control tracking performance for the dynamic walking gait (\cref{sec:BipedSimulation}) and a highly-dynamic forward jump (\cref{sec:HighlySimulation}).

\subsubsection{Dynamic Walking}
To begin with, we analyze the motion tracking for dynamic walking gait. \cref{fig:walkDynamic_pybulletTracking} shows the tracking performance of the joint level control architecture. The reference trajectories from \gls{OC} are visualized as solid lines, while the resulting trajectories of the real-time physics simulation are shown as dotted lines. It can be seen that the controller follows the optimal trajectories well. Small deviations can be seen around two seconds, which accounts for the lift off phase for the second robot step. This effect can be explained by the apparent abrupt change in the joint space but is found to be marginal for the overall tracking performance. 
\begin{figure}
\centering	
\includegraphics[width=.85\textwidth]{fig/walkDynamic/pybullet/pybulletTracking}
\caption[Control tracking performance on joint level for dynamic walking]{Control tracking performance on joint level for dynamic walking. The results indicate that the control architecture follows the reference trajectories with sufficiently good accuracy.}
\label{fig:walkDynamic_pybulletTracking}
\end{figure}
As introduced previously, the control architecture is solely based on joint space position and velocity level. Although the tracking performance is good, this may not prove for the stability of the motions. Similar to the definition of robot tasks, the evaluation of the motions should be pursued in task space. To this end, \cref{fig:walkDynamic_pybulletBase} monitors the according motion of the floating base. As can be seen, the floating base deviates about $\pm$ 10 mm in x and y-direction as well as $+$ 5 mm in z-direction. 
\begin{figure}
\centering	
\includegraphics[width=.85\textwidth]{fig/walkDynamic/pybullet/pybulletBase}
\caption[Motion of the floating base for dynamic walking]{Motion of the floating base resulting from joint space control for the dynamic walking gait. The results show a deviation of the floating base about $\pm$ 10 mm in x and y-direction , as well as $+$ 5 mm in z-direction.}
\label{fig:walkDynamic_pybulletBase}
\end{figure}

In this context, it is important to notice that no controller tracks these task space quantities. Instead, they are merely the result of the joint space control performance. The largest deviations in task space occur during the first step and the according impact in the first two seconds of the motions. It becomes evident that these task space errors do not correlate to the peaks discussed on joint space level. Furthermore, one can observe oscillations in the stabilization phase at the end of the motion. This effect can be attributed to the fact that the systems slightly starts to swing after the second impact. This behavior could be compensated e.g. with a dedicated control in task space instead of joint space. 

Overall, it can be stated that the dynamic walking motion can be successfully stabilized by the proposed control approach. 

\subsubsection{Forward Jumping}
To identify the limits of the control approach, we now study the online stabilization for a highly-dynamic forward jump (\cref{sec:HighlySimulation}). In addition to the previously discussed dynamic walking, the jumping task introduces new challenges for the controller in terms of speed and robustness that will be investigated in the following.  

\cref{fig:jumpForward_pybulletTracking} shows the tracking performance of the joint level control architecture for the forward jump. In contrast to the case of dynamic walking, the joint space controller reveals larger tracking deviations. This is especially true for the most dynamic part of the motion, namely the acceleration of the base and finally the takeoff around 300-400 ms. Large deviations can be especially seen for body pitch, and the knee joints, which turned out to be crucial for highly-dynamic movements. 

\begin{figure}[h!]
\centering	
\includegraphics[width=.85\textwidth]{fig/jumpForward/pybullet/pybulletTracking}
\caption[Control tracking performance on joint level for a forward jump]{Control tracking performance on joint level for a forward jump. The results indicate that the control architecture also is suitable to track the reference trajectories of highly-dynamic movements with sufficiently good accuracy.}
\label{fig:jumpForward_pybulletTracking}
\end{figure}

In addition to these findings, \cref{fig:jumpForward_pybulletBase} monitors the motion of the (uncontrolled) floating base resulting from the joint space control performance. As becomes evident, the deviations in task space are also much higher compared to the dynamics walking case. While the height of the floating base is reasonable, the x-position shows tracking errors of about $\pm$ 5 cm. Errors of this magnitude inevitably lead to instability of the movement to be performed. In this case, the strong deviation in task space causes the robot to tilt around the rear edge of the foot after the touchdown. 

\begin{figure}[h!]
\centering	
\includegraphics[width=.85\textwidth]{fig/jumpForward/pybullet/pybulletBase}
\caption[Motion of the floating base for a forward jump]{Motion of the floating base for a forward jump.}
\label{fig:jumpForward_pybulletBase}
\end{figure}

Consequently, it turned out that a mere control on a joint space basis is not sufficient to track highly dynamic movements but is indeed appropriate to stabilize a dynamic walking gait in real-time. 


\section{Validation in Real-World Experiments}\label{sec:OnlineExperiments}
This section investigates the stability of the planned motions in real-world experiments on the full-size humanoid robot RH5 (see \cref{subsec:RH5}). Analogously to the validation in PyBullet, the goal is to track the \gls{OC} trajectories with a joint space online stabilization on the real system. 

A total of four experiments are conducted. The primary goal is to examine whether the presented control architecture is sufficient to follow the static and dynamic walking trajectories (see \crefrange{c4}{c5}). For a dedicated root cause analysis, we test two additional motions on the real-system, namely one-legged balancing and a sequence of dynamic squat movements. The respective goal is to identify potential model deviations and to study the capabilities of the real-time controller.

Based on the simulation results from the last section, highly dynamic movements are not evaluated on the real system in the context of this thesis. As discussed in the last section, these motions definitely require more advanced control algorithms in task space due to the dynamic nature of the movement. 

The rest of this section is structured as follows. First, an overview of the experimental setup with the involved components is presented. Following up, the experiments are presented and discussed with respect to the tracking performance and the stability of the motion.

\subsection{Pipeline}\label{subsec:Pipeline}
The foundation for the experiments are the motion data generated offline with the proposed whole-body \gls{TO}. The planned motions are then tracked in real-time with a joint space online stabilization on the real system. An overview of the experimental pipeline is given in \cref{img:experimentalSetup}. In the following, details on the involved components are provided. 
\begin{figure}
\centering	
\includegraphics[width=1\textwidth]{img/experimentalSetup2}
\caption[Overview about the experimental pipeline]{Overview about the experimental pipeline. The offline planned motions are tracked in real-time with a joint space online stabilization on the real system.}
\label{img:experimentalSetup}
\end{figure}
%Trajectory File
The presented motion planning approach computes inherently balanced motions that are concisely captured in an appropriate file. This trajectory file contains the optimal state trajectories $\myM{X}^*=[\bq^*,\bv^*]$, optimal control inputs $\myM{U}^*$ and the resulting contact wrenches $\myM{F}_{\text{ ext}}^*$ acting on the feet. In order to minimize the computational effort in the real-time loop, the file already encompasses data discretized to the desired frequency of 1 kHz. As in the PyBullet validation, the trajectories are interpolated using cubic splines in order to ensure smoothness and differentiability. 

%HyRoDyn motivation
As introduced in \cref{subsec:RH5}, the novel RH5 humanoid robot contains multiple parallel mechanisms in order to achieve a high dynamic performance, superior stiffness and payload-to-weight ratio. This leads to the presence of various closed loops and hence a series-parallel hybrid robotic system, which is difficult to model and control. In most multi-body dynamics libraries, e.g. RBDL \cite{felis2017rbdl} and OpenSim \cite{delp2007opensim}, these loop closure constraints are solved numerically. HyRoDyn (Hybrid Robot Dynamics), is a recently presented modular software framework for solving the kinematics and dynamics of these type of series-parallel hybrid robots analytically, leading to improved accuracy and computational performance \cite{kumar2020analytical}.

%Abstract vs Full Model
The planned motions are computed based on a tree type robot model. For dynamic real-time control, this simplified model turns out to be sufficient, although the accuracy is reduced \cite{kumar2019model}. Nevertheless, the problem remains on transforming the results from the independent joint space, to the actuation space. In the context of this thesis, HyRoDyn is used to map the trajectories generated for the serialized robot model to compute the forces of the respective linear actuators.

%Framwork Consistencies
Consistency of the frameworks involved in the motion planning and control pipeline is an indispensable prerequisite for the following experiments. Hence, \cref{app:Consistency} provides a brief verification of the notability consistency between HyRoDyn and Pinocchio \cite{carpentier2019pinocchio}, which is used inside Crocoddyl for computation of robot dynamics and their analytical derivatives. 

%Control+Rock
Low-level actuator controllers are utilized to compensate deviations from the reference trajectories. Analogously to the PyBullet validation, this control approach uses a cascaded feedback of position, velocity and an additional current control loop. 
%HyRoDyn, as well as the low-level control of the RH5 robot is implemented using the Robot Construction Kit (Rock) middleware \cite{joyeux2013rock}, which is based on the Orocos Real Time Toolkit and forms the software backbone of the RH5 humanoid robot. 

\subsection{Experiments}
We conduct a total of four experiments, gradually incorporating a new level of difficulty (see \cref{tab:experiments}), that are introduced in the following.
\begin{table}[t]
\centering
\caption[Overview about the difficulty of the conducted experiments]{Overview about the difficulty of the conducted experiments. Three out of four planned motions could be successfully stabilized by the controller on the real system.}
\begin{tabular}{lcccc}
\hline
& Balancing & Static Walk & Fast Squats & Dynamic Walk\\ \hline
Surface contacts & \greencheckmark  & \greencheckmark & \greencheckmark & \greencheckmark \\
Base motion & \greencheckmark  & \greencheckmark & \greencheckmark & \greencheckmark \\
Swing foot motion & \greencheckmark  & \greencheckmark & \redxmark & \greencheckmark \\
Step sequence & \redxmark  & \greencheckmark & \redxmark & \greencheckmark \\
Impacts & \redxmark  & \greencheckmark & \redxmark & \greencheckmark \\
Dynamic forces & \redxmark  & \redxmark & \greencheckmark & \greencheckmark \\
%Non-flat terrain & \redxmark  & \redxmark & \redxmark & \redxmark \\
Flight-phases & \redxmark  & \redxmark & \redxmark & \redxmark \\
\hline
\quad\quad Success & \greencheckmark  & (\greencheckmark) & \greencheckmark & \redxmark \\ \hline
\end{tabular}
\label{tab:experiments}
\end{table}
\subsubsection{Experiment I: One-Leg Balancing}
The goal of this first experiment is to test the ability of the control architecture to track a slow balancing task. The quasi-static motion consists of five locomotion phases as visualized in \cref{exp:balancingSnaps}. From a (a) stable pose in \gls{DS}, (b) shifting the \gls{CoM} above the \gls{LF} (c) lifting the \gls{RF} slightly up (d) and down and (e) returning to the initial pose.
\begin{figure}
\begin{subfigure}{.2\textwidth}
	\includegraphics[width=.95\linewidth]{experiments/balancing/snaps/1}
	\caption{}
	\end{subfigure}%
\begin{subfigure}{.2\textwidth}
	\includegraphics[width=.95\linewidth]{experiments/balancing/snaps/2}
	\caption{}
\end{subfigure}%
\begin{subfigure}{.2\textwidth}
	\includegraphics[width=.95\linewidth]{experiments/balancing/snaps/3}
	\caption{}
	\end{subfigure}%
\begin{subfigure}{.2\textwidth}
	\includegraphics[width=.95\linewidth]{experiments/balancing/snaps/4}
	\caption{}
\end{subfigure}%
\begin{subfigure}{.2\textwidth}
	\includegraphics[width=.95\linewidth]{experiments/balancing/snaps/5}
	\caption{}
	\end{subfigure}%
\caption[Experiment I: One-leg balancing]{Experiment I: One-leg balancing from (a) an initial pose, (b) \gls{CoM} shift above the \gls{LF}, (c) lifting the \gls{RF} up and (d) down and (e) recovering to the initial pose. \href{https://github.com/julesser/ma-thesis-experimental-results/blob/master/balancing/balancing.mp4}{[Video]}} 
\label{exp:balancingSnaps}
\end{figure} 

\subsubsection{Experiment II: Static Walking}
The second experiment deals with a stabilization of the static walking pattern discussed in \cref{sec:BipedSimulation} (see \cref{exp:staticWalkingSnaps}). The objective of this test is to analyze the effect of more difficult swing-leg motions, a step sequence of two steps and the effect of impacts. 
\begin{figure}
\begin{subfigure}{.2\textwidth}
	\includegraphics[width=.95\linewidth]{experiments/walkStatic/snaps/1}
	\caption{}
	\end{subfigure}%
\begin{subfigure}{.2\textwidth}
	\includegraphics[width=.95\linewidth]{experiments/walkStatic/snaps/2}
	\caption{}
\end{subfigure}%
\begin{subfigure}{.2\textwidth}
	\includegraphics[width=.95\linewidth]{experiments/walkStatic/snaps/3}
	\caption{}
	\end{subfigure}%
\begin{subfigure}{.2\textwidth}
	\includegraphics[width=.95\linewidth]{experiments/walkStatic/snaps/4}
	\caption{}
\end{subfigure}%
\begin{subfigure}{.2\textwidth}
	\includegraphics[width=.95\linewidth]{experiments/walkStatic/snaps/5}
	\caption{}
	\end{subfigure}%
\caption[Experiment II: Static walking]{Experiment II: Static walking from (a) an initial pose, (b) \gls{CoM} shift above the \gls{LF}, (c,d) performing a right step and (e) shifting the \gls{CoM} to the center of the \gls{SP}. \href{https://github.com/julesser/ma-thesis-experimental-results/tree/master/walkStatic}{[Video]}}
\label{exp:staticWalkingSnaps}
\end{figure} 

\subsubsection{Experiment III: Fast Squats}
The objective of this third experiment is to evaluate the tracking performance in the context of a dynamic motion. In contrast to the first two motions, the fast squatting experiment (see \cref{exp:squatSnaps}) involves dynamic forces acting on the robot resulting from a fast vertical base movement in the range of 15 cm within two seconds. Analogously to the multiple forward jumps (see \cref{sec:BipedSimulation}), the motion is composed of a sequence of three \gls{OC} problems that are solved sequentially (see \cref{eqn:optimizationProblemSequence}). This experiment can be seen as preliminary test for dynamic walking following in the next experiment.   
\begin{figure}
\begin{subfigure}{.14\textwidth}
	\includegraphics[width=.95\linewidth]{experiments/squats/snaps/1}
	\caption{}
	\end{subfigure}%
\begin{subfigure}{.14\textwidth}
	\includegraphics[width=.95\linewidth]{experiments/squats/snaps/2}
	\caption{}
\end{subfigure}%
\begin{subfigure}{.14\textwidth}
	\includegraphics[width=.95\linewidth]{experiments/squats/snaps/3}
	\caption{}
	\end{subfigure}%
\begin{subfigure}{.14\textwidth}
	\includegraphics[width=.95\linewidth]{experiments/squats/snaps/4}
	\caption{}
\end{subfigure}%
\begin{subfigure}{.14\textwidth}
	\includegraphics[width=.95\linewidth]{experiments/squats/snaps/5}
	\caption{}
	\end{subfigure}%
\begin{subfigure}{.14\textwidth}
	\includegraphics[width=.95\linewidth]{experiments/squats/snaps/6}
	\caption{}
	\end{subfigure}%
\begin{subfigure}{.14\textwidth}
	\includegraphics[width=.95\linewidth]{experiments/squats/snaps/7}
	\caption{}
	\end{subfigure}%
\caption[Experiment III: Sequence of fast squats]{Experiment III: Sequence of fast squats from (a) an initial pose over, (b,d,f) descending the \gls{CoM} by 15 cm and (c,e,g) recovering to the initial pose. \href{https://github.com/julesser/ma-thesis-experimental-results/blob/master/squats/squat15cm_2s.mp4}{[Video]}}
\label{exp:squatSnaps}
\end{figure} 

\subsubsection{Experiment IV: Dynamic Walking}
This experiment investigates the capabilities of the control approach to follow the dynamic walking gait trajectories presented in \cref{sec:BipedSimulation} (see \cref{exp:dynamicWalkingSnaps}). From a perspective of complexity, it combines both the challenges of static walking, namely dedicated swing-leg motions, step sequence and impacts, with the difficulty of dynamic forces acting on the robot as explored with the fast squats experiment.
\begin{figure}
\begin{subfigure}{.2\textwidth}
	\includegraphics[width=.95\linewidth]{experiments/walkDynamic/snaps/1}
	\caption{}
	\end{subfigure}%
\begin{subfigure}{.2\textwidth}
	\includegraphics[width=.95\linewidth]{experiments/walkDynamic/snaps/2}
	\caption{}
\end{subfigure}%
\begin{subfigure}{.2\textwidth}
	\includegraphics[width=.95\linewidth]{experiments/walkDynamic/snaps/3}
	\caption{}
	\end{subfigure}%
\begin{subfigure}{.2\textwidth}
	\includegraphics[width=.95\linewidth]{experiments/walkDynamic/snaps/4}
	\caption{}
\end{subfigure}%
\begin{subfigure}{.2\textwidth}
	\includegraphics[width=.95\linewidth]{experiments/walkDynamic/snaps/5}
	\caption{}
	\end{subfigure}%
\caption[Experiment IV: Dynamic walking]{Experiment IV: Dynamic walking from (a) an initial pose, (b,c) performing a right step and (e,f) a successive left step. \href{https://github.com/julesser/ma-thesis-experimental-results/blob/master/walkDynamic/walkDynamicFail.mp4}{[Video]}}
\label{exp:dynamicWalkingSnaps}
\end{figure} 
eo
\subsection{Discussion}
The following section analyzes the experimental results, highlighting limits of the utilized control architecture.
\subsubsection{Stability of the Motions}
All in all, three out of four planned motions could be successfully stabilized by the  controller on the real system (see \cref{tab:experiments}). Both the one-leg balancing experiment (I) and the fast squats experiment (III) were stabilized with good accuracy. For experiment II, static walking, only one step could be stabilized at a time. Furthermore, it could be shown that dynamic walking (IV) cannot be stabilized with the present control approach. Possible root causes are outlined in the following. 
\subsubsection{Tracking Performance}
As detailed in \cref{subsec:Pipeline}, the online stabilizer solely works in joint space. The experimental results indicate that the overall joint space tracking for all four experiments is satisfying. 
\cref{exp:BalancingTracking} exemplary shows the tracking performance for the one-leg balancing experiment. It becomes evident that the control architecture allows following the computed reference trajectory in actuator space (a,b) closely. Furthermore, also the resulting tracking quality in the virtual joints (c,d) is satisfying. Similar conclusions hold for the other experiments. Even during critical impulse phases the joint space tracking shows no abnormalities. This precise tracking is achieved with high-gain joint space control, which allows a quick compensation of position differences that comes at the cost of lost compliance in the joints.
\begin{figure}[h!]
\begin{subfigure}{.5\textwidth}
	\includegraphics[page=1, width=.95\linewidth]{experiments/balancing/ankleTrack}
	\caption{LLAnkleAct1}	
	\end{subfigure}%
\begin{subfigure}{.5\textwidth}
	\includegraphics[page=2, width=.95\linewidth]{experiments/balancing/ankleTrack}
	\caption{LLAnkleAct2}	
\end{subfigure}%

\begin{subfigure}{.5\textwidth}
	\includegraphics[page=3, width=.95\linewidth]{experiments/balancing/ankleTrack}
	\caption{LLAnkleRoll}	
	\end{subfigure}%
\begin{subfigure}{.5\textwidth}
	\includegraphics[page=4, width=.95\linewidth]{experiments/balancing/ankleTrack}
	\caption{LLAnklePitch}	
\end{subfigure}%
\caption[Tracking performance for the one-leg balancing experiment]{Tracking performance for the one-leg balancing experiment in actuator space (a,b) and independent joint space (c,d). The reference position from \gls{OC} in orange and the measured position in blue are plotted against the time. The results indicate that the control architecture tracks the reference trajectories with sufficient accuracy.}
\label{exp:BalancingTracking}
\end{figure} 
\subsubsection{Handling Impulses}
The impact phase turned out to be the main problem for the walking experiments. This is reasonable since the utilized control approach only compensates for errors in joint space, while errors in task space can arise quickly and are not compensated. Deviations of the reference base motion only effect the static stability if the \gls{FCoM} leaves the support polygon. Differences from the foot position instead may cause instability with a respective error in the range of a few millimeters. The same holds for errors in the foot orientation. Consider for example the case in which the robot only touches the ground with the rear part of the swing foot's sole. Instead of the assumed surface contact between the foot and the ground, a line contact appears that causes the robot to fall in a combination with high-gain position control. 
\subsubsection{Model Discrepancies}
The quasi-static experiments provide a good basis of analyzing the model accuracy. Due to the small \gls{SP}, model deviations inevitably lead to instability during the single support phases. In the balancing and static walking experiments the motions could be stabilized if, and only if, the \gls{FCoM} was adjusted by about $\pm$ 7.5 mm. The found model errors in the mass distribution, possibly combined with with errors in inertia distribution, may also pose part of the problem for walking gait stabilization. In order to compensate for these differences, an extensive system identification could be performed to check the dynamic model against the real system.
\subsubsection{Mechanical Deficiencies}
Beneath the identified limits of the control approach and the modeling errors, also mechanical weaknesses of the RH5 humanoid may contribute to the sim-to-real gap. In the scope of the experiments, several issues have been revealed, namely joint play, calibration errors and undesired structural flexibility in the ankle of the foot. Furthermore, the rigid contact points on the sole of the feet are found to be not the ideal solution, since small deviations can quickly lead to instability of the whole system. Instead, it may be advantageous for dynamic movements of the humanoid robot to use a planar sole, providing adequate damping properties. 



























%-----------------------------------------------------------------------------%
%                                                                             %
%    K A P I T E L   7                                                       %
%                                                                             %
%-----------------------------------------------------------------------------%

\chapter{Conclusion and Outlook}\label{c7}

\section{Thesis Summary}
This master thesis was motivated by the generation of physically consistent, efficient motions plans for legged robots. 

The premise of investigation was that whole-body planning leads to more efficient motions than a simple program such as an \gls{IK} solver. Therefore, the core algorithm of the proposed motion planning approach is a recently presented \gls{DDP}-based whole-body \gls{TO}. Building upon this, a generic method for constraining DDP-like solvers is presented in order to generate dynamically balanced motions. The results are integrated into the recently presented open-source framework Crocoddyl.

Following this, we investigated the \gls{CoP}-based contact stability of the  proposed motion planning approach for a wide range of motions with the biologically inspired RH5 humanoid robot. It has been shown that the resulting motion plans for both, dynamic bipedal walking and various jumping tasks, are inherently balanced. Additionally, the analysis of highly-dynamic movements allowed the derivation of useful guidelines for future design iterations of the humanoid robot.

Although the focus of this thesis was on motion planning, we evaluated the feasibility of the generated trajectories with a simple online stabilizer. It has been demonstrated in a real-time physics simulator that the motion plans can be stabilized by a simple control architecture solely based on joint-space position control. Furthermore, it could be shown that for real-world experiments, a control in task space is indispensable in order to compensate deviations between model and reality. 

The final result of this thesis is an efficient motion planning approach that produces inherently balanced motions. This algorithm efficiently generates highly-dynamic movements with flight-phases for various legged systems. 

\section{Future Directions}
%DDP is great
We see large potential in using \gls{DDP}-based whole-body \gls{TO} to generate motions for legged systems. Motion planning based on numerical optimization clearly reduces the amount of hand-crafted components and allows the specification of high-level tasks directly in the operational space of the robot. This might become even more important as legged robots tackle more difficult terrain that require higher dynamic motions.
  
%Algorithmic Improvements
From an algorithmic perspective, the formulation can still be improved in a number of ways. The goal for the algorithm is to efficiently and accurately generate physically compliant motions. Promising ways of simultaneously improving both measures are seen by embedding inequality constraints directly as strict bounds inside the \gls{DDP} algorithm, instead of forcing it by penalization.

%Control Approaches: MPC + Online Stabilization
From a control perspective, two successive steps are of particular interest. First, it is considered worth working on an improved online stabilization. As discussed in the last chapter, operational space control is inevitable in order to compensate for modeling errors directly in tasks space. Following up, it would be interesting to embed the motion planning inside a \gls{MPC} formulation. If the re-planning is fast enough, the robot is enabled to act more robustly to unpredicted situations. 

%Long-term vision
This direction of future research, namely improving the motion-planning algorithm and embedding it in an \gls{MPC} formulation on a real system, seems promising. Intelligently combining optimization-based planning and control may be the key for robots to interact with both environment and humans in a more natural, dynamic and autonomous way. With some work in this direction, we might soon find legged robots crossing our daily path or intuitively collaborating with them when assembling infrastructures and exploring foreign planets. 

%-----Anhang-----
%-----------------------------------------------------------------------------%
%                                                                             %
%    Anh�nge					                                             					  %
%                                                                             %
%-----------------------------------------------------------------------------%

\appendix
\chapter{Appendix}

\subsection{Carlos Talk: Essentials}
In the end we want to solve a bilevel (nested) optimization
\begin{equation} 
\begin{split}
\myM{X}^*,\myM{U}^* = \arg\min_{\myM{X},\myM{U}} \sum_{k=0}^{N-1} task(x_k,u_k) \\
x_k=\arg\min physics(x_k,u_k), \\
s.t. \quad constraints(x_k,u_k)
\end{split}
\end{equation}

Which more formally looks like
\begin{equation*} 
\begin{split}
\myM{X}^*,\myM{U}^*=
\begin{Bmatrix} \myM{x}^*_0,\cdots,\myM{x}^*_N \\
      \myM{u}^*_0,\cdots,\myM{u}^*_N
\end{Bmatrix} =
\arg\min_{\myM{X},\myM{U}} l_N(x_N)+\sum_{k=0}^{N-1} \int_{t_k}^{t_k+\Delta t} l(\myM{x},\myM{u})dt \\
s.t.\quad\myM{\dot{v}}, \myM{\lambda}=\arg\min_{\myM{\dot{v}}, \myM{\lambda}}||\myM{\dot{v}}-\myM{\dot{v}}_{free}||_M, \\
\myM{x}\in \mathcal{X}, \myM{u}\in \mathcal{U}
\end{split}
\end{equation*}

KKT Matrix:
\begin{equation*} 
\begin{split}
\myM{X}^*,\myM{U}^* = \arg\min_{\myM{X},\myM{U}} \sum_{k=0}^{N-1} task(x_k,u_k) \\
KKT-Dynamics(x_k,u_k) 
\end{split}
\end{equation*}

Multi-contact dynamics as holonomic constraints:
\begin{equation*} 
\left[\begin{matrix}
  \myM{\dot{v}} \\
  -\myM{\lambda} \\
\end{matrix}\right] 
=
\left[\begin{matrix}
 \myM{M} & \myM{J}^{\top}_c \\
 {\myM{J}_{c}} & \myM{0} \\
\end{matrix}\right]^{-1}
\left[\begin{matrix}
 \myM{\tau}_b \\ -\myM{a_0}
\end{matrix}\right]
\end{equation*}


\section{Crocoddyl: Contact RObot COntrol by Differential DYnamic programming Library (Wiki Home)}
\subsection{Welcome to Crocoddyl}
Crocoddyl is an optimal control library for robot control under contact sequence. Its solver is based on an efficient Differential Dynamic Programming (DDP) algorithm. Crocoddyl computes optimal trajectories along to optimal feedback gains. It uses Pinocchio for fast computation of robot dynamics and its analytical derivatives.
Crocoddyl is focused on multi-contact optimal control problem (MCOP) which as the form:
$$\myM{X}^*,\myM{U}^*= \begin{Bmatrix} \myM{x}^*_0,\cdots,\myM{x}^*_N \\ \myM{u}^*_0,\cdots,\myM{u}^*_N \end{Bmatrix} = \arg\min_{\myM{X},\myM{U}} \sum_{k=1}^N \int_{t_k}^{t_k+\Delta t} l(\myM{x},\myM{u})dt$$
subject to
$$\myM{\dot{x}} = \myM{f}(\myM{x},\myM{u}),$$
$$\myM{x}\in\mathcal{X}, \myM{u}\in\mathcal{U}, \boldsymbol{\lambda}\in\mathcal{K}.$$
where
\begin{itemize}
\item the state $\myM{x}=(\myM{q},\myM{v})$ lies in a manifold, e.g. Lie manifold $\myM{q}\in SE(3)\times \myM{R}^{n_j}$,
\item the system has underactuacted dynamics, i.e. $\myM{u}=(\myM{0},\boldsymbol{\tau})$,
\item $\mathcal{X}$, $\mathcal{U}$ are the state and control admissible sets, and
\item $\mathcal{K}$ represents the contact constraints.
\end{itemize}

Note that $\boldsymbol{\lambda}=\myM{g}(\myM{x},\myM{u})$ denotes the contact force, and is dependent on the state and control.

Let's start by understanding the concept behind crocoddyl design.

\subsection{Action Models}
In crocoddyl, an action model combines dynamics and cost models. Each node, in our optimal control problem, is described through an action model. Every time that we want describe a problem, we need to provide ways of computing the dynamics, cost functions and their derivatives. All these is described inside the action model.
To understand the mathematical aspects behind an action model, let's first get a locally linearize version of our optimal control problem as:
$$\myM{X}^*(\myM{x}_0),\myM{U}^*(\myM{x}_0) = \arg\min_{\myM{X},\myM{U}} = cost_T(\delta\myM{x}_N) + \sum_{k=1}^N cost_t(\delta\myM{x}_k, \delta\myM{u}_k)$$
subject to
$$dynamics(\delta\myM{x}_{k+1},\delta\myM{x}_k,\delta\myM{u}_k)=\myM{0},$$
where
$$cost_T(\delta\myM{x}_k) = \frac{1}{2} \begin{bmatrix} 1 \\ \delta\myM{x}_k \end{bmatrix}^\top \begin{bmatrix} 0 & \myM{l_x}^\top_k \\ \myM{l_x}_k & \myM{l_{xx}}_k \end{bmatrix} \begin{bmatrix} 1 \\ \delta\myM{x}_k \end{bmatrix},$$     

$$cost_t(\delta\myM{x}_k,\delta\myM{u}_k) = \frac{1}{2} \begin{bmatrix} 1 \\ \delta\myM{x}_k \\ \delta\myM{u}_k \end{bmatrix}^\top \begin{bmatrix} 0 & \myM{l_x}^\top_k & \myM{l_u}^\top_k\\ \myM{l_x}_k & \myM{l_{xx}}_k & \myM{l_{ux}}^\top_k\\ \myM{l_u}_k & \myM{l_{ux}}_k & \myM{l_{uu}}_k \end{bmatrix} \begin{bmatrix} 1 \\ \delta\myM{x}_k \\ \delta\myM{u}_k \end{bmatrix}$$

$$dynamics(\delta\myM{x}_{k+1},\delta\myM{x}_k,\delta\myM{u}_k) = \delta\myM{x}_{k+1} - (\myM{f_x}_k\delta\myM{x}_k + \myM{f_u}_k\delta\myM{u}_k)$$

\subsubsection{Notes}
\begin{itemize}
\item An action model describes the dynamics and cost functions for a node in our optimal control problem.
\item Action models lie in the discrete time space.
\item For debugging and prototyping, we have also implemented NumDiff abstractions. These computations depend only in the defining of the dynamics equation and cost functions. However to asses efficiency, crocoddyl uses analytical derivatives computed from Pinocchio.
\end{itemize}

\subsubsection{Differential and Integrated Action Models}

It's often convenient to implement action models in continuous time. In crocoddyl, this continuous-time action models are called Differential Action Model (DAM). And together with predefined Integrated Action Models (IAM), it possible to retrieve the time-discrete action model needed by the solver.
At the moment, we have the following integration rules:
\begin{itemize}
\item simpletic Euler and
\item Runge-Kutta 4.
\end{itemize}

\subsubsection{Add On from Introduction.jpnb}
Optimal control solvers often need to compute a quadratic approximation of the action model (as previously described); this provides a search direction (computeDirection). Then it's needed to try the step along this direction (tryStep).

Typically calc and calcDiff do the precomputations that are required before computeDirection and tryStep respectively (inside the solver). These functions update the information of:
\begin{itemize}
\item \textbf{calc}: update the next state and its cost value
 $$\delta\myM{\dot{x}}_{k+1} = \myM{f}(\delta\myM{x}_k,\myM{u}_k)$$
\item \textbf{calcDiff}: update the derivatives of the dynamics and cost (quadratic approximation)
 $$\myM{f_x}, \myM{f_u} \hspace{1em} (dynamics)$$
 $$\myM{l_x}, \myM{l_u}, \myM{l_{xx}}, \myM{l_{ux}}, \myM{l_{uu}} \hspace{1em} (cost)$$
\end{itemize}


\subsection{State and its Integrate and Difference Rules}
General speaking, the system's state can lie in a manifold $M$ where the state rate of change lies in its tangent space $T_\mathbf{x} M$. There are few operators that needs to be defined for different routines inside our solvers:
$$\myM{x}_{k+1} = integrate(\myM{x}_k,\delta\myM{x}_k) = \myM{x}_k \oplus \delta\myM{x}_k$$
$$\delta\myM{x}_k = difference(\myM{x}_{k+1},\myM{x}_k) = \myM{x}_{k+1} \ominus \myM{x}_k$$
where $\mathbf{x}\in M$ and $\delta\mathbf{x}\in T_\mathbf{x} M$.
And we also need to defined the Jacobians of these operators with respect to the first and second arguments:
$$\frac{\partial \myM{x}\oplus\delta\myM{x}}{\partial \myM{x}}, \frac{\partial \myM{x}\oplus\delta\myM{x}}{\partial\delta\myM{x}} =Jintegrante(\myM{x},\delta\myM{x})$$
$$\frac{\partial\myM{x}_2\ominus\myM{x}_2}{\partial \myM{x}_1}, \frac{\partial \myM{x}_2\ominus\myM{x}_1}{\partial\myM{x}_1} =Jdifference(\myM{x}_2,\myM{x}_1)$$
For instance, a state that lies in the Euclidean space will the typical operators:
$$integrate(\myM{x},\delta\myM{x}) = \myM{x} + \delta\myM{x}$$
$$difference(\myM{x}_2,\myM{x}_1) = \myM{x}_2 - \myM{x}_1$$
$$Jintegrate(\cdot,\cdot) = Jdifference(\cdot,\cdot) = \myM{I}$$
All these functions are encapsulate inside the State class. For Pinocchio models, we have implemented the StateMultibody class which can be used for any robot model.

\section{Crocoddyl Wiki: Differential Action Model for Floating in Contact Systems (DAMFIC)}
\subsection{System Dynamics}
As you might know, a differential action model describes the systems dynamics and cost function in continuous-time. For multi-contact locomotion, we account for the rigid contact by applying the Gauss principle over holonomic constraints in a set of predefined contact placements, i.e.:
$$\begin{aligned} & \dot{\myM{v}} = \underset{\myM{a}}{\arg\min} & & \frac{1}{2}\,\|\dot{\myM{v}}-\dot{\myM{v}}_{free}\|_{\myM{M}} \\ & \textrm{subject to} & & \myM{J}_{c} \dot{\myM{v}} + \dot{\myM{J}}_c \myM{v} = \myM{0}, \end{aligned}$$

This is equality-constrained quadratic problem with an analytical solution of the form:
$$\left[\begin{matrix}\myM{M} & \myM{J}^{\top}_c \\{\myM{J}_{c}} & \myM{0}\end{matrix}\right] \left[\begin{matrix} \dot{\myM{v}} \\ -\boldsymbol{\lambda} \end{matrix}\right] = \left[\begin{matrix} \boldsymbol{\tau}_b \\ -\dot{\myM{J}}_c \myM{v}\end{matrix}\right]$$
in which
$$(\dot{\myM{v}},\boldsymbol{\lambda})\in(\myM{R}^{nv},\myM{R}^{nf})$$
are the primal and dual solutions,
$$\myM{M}\in\myM{R}^{nv\times nv}$$ is formally the metric tensor over the configuration manifold $\myM{q}\in\myM{R}^{nq}$,
$$\myM{J}_{c}= \begin{bmatrix} \myM{J}_{c_1} & \cdots & \myM{J}_{c_f}\end{bmatrix}\in\myM{R}^{nf\times nv}$$ is a stack of $f$ contact Jacobians,
$\boldsymbol{\tau}_b=\myM{S}\boldsymbol{\tau}-\myM{b}\in\myM{R}^{nv}$ is the force-bias vector that accounts for the control $\boldsymbol{\tau}\in\myM{R}^{nu}$, the Coriolis and gravitational effects $\myM{b}$, and $\myM{S}$ is the selection matrix of the actuated joint coordinates, and
$nq$, $nv$, $nu$ and $nf$ are the number of coordinates used to describe the configuration manifold, its tangent-space dimension, control commands and contact forces, respectively.

And this equality-constrained forward dynamics can be formulated using state space representation, i.e.:
$$\dot{\myM{x}}=\myM{f(x,u)}$$
where $\myM{x}=(\myM{q,v})\in\myM{R}^{nq+nv}$ and $\myM{u}=\boldsymbol{\tau}\in\myM{R}^{nu}$ are the state and control vectors, respectively. Note that $\dot{\myM{x}}$ lies in the tangent-space of $\myM{x}$, and their dimension are not the same.

\subsection{Add On from Introduction.jpnb}
\subsection{Solving the Optimal Control Problem}
Our optimal control solver interacts with a defined ShootingProblem. A \textbf{shooting problem} represents a \textbf{stack of action models} in which an action model defines a specific node along the OC problem.

First we need to create an action model from DifferentialFwdDynamics. We use it for building terminal and running action models. In this example, we employ an simpletic Euler integration rule.

Next we define the set of cost functions for this problem. One could formulate
\begin{itemize}
\item Running costs (related to individual states)
\item Terminal costs (related to the final state)
\end{itemize}
in order to penalize, for example, the state error, control error, or end-effector pose error. 

Onces we have defined our shooting problem, we create a DDP solver object and pass some callback functions for analysing its performance.
\subsubsection{Application to Bipedal Walking}
In crocoddyl, we can describe the multi-contact dynamics through holonomic constraints for the support legs. From the Gauss principle, we have derived the model as:
$$
\left[\begin{matrix}
 \myM{M} & \myM{J}^{\top}_c \\
 {\myM{J}_{c}} & \myM{0} \\
\end{matrix}\right]
\left[\begin{matrix}
 \dot{\myM{v}} \\ -\boldsymbol{\lambda}
\end{matrix}\right]
 = 
\left[\begin{matrix}
  \boldsymbol{\tau} - \myM{h} \\
  -\dot{\myM{J}}_c \myM{v} \\
\end{matrix}\right]$$.

This DAM is defined in "DifferentialActionModelFloatingInContact" class.
Given a predefined contact sequence and timings, we build per each phase a specific multi-contact dynamics. Indeed we need to describe \textbf{multi-phase optimal control problem}. One can formulate the multi-contact optimal control problem (MCOP) as follows:

$$\myM{X}^*,\myM{U}^*=
\begin{Bmatrix} \myM{x}^*_0,\cdots,\myM{x}^*_N \\
				  \myM{u}^*_0,\cdots,\myM{u}^*_N
\end{Bmatrix} =
\arg\min_{\myM{X},\myM{U}} \sum_{p=0}^P \sum_{k=1}^{N(p)} \int_{t_k}^{t_k+\Delta t} l_p(\myM{x},\myM{u})dt$$
subject to
$$ \myM{\dot{x}} = \myM{f}_p(\myM{x},\myM{u}), \text{for } t \in [\tau_p,\tau_{p+1}]$$

$$ \myM{g}(\myM{v}^{p+1},\myM{v}^p) = \myM{0}$$

$$ \myM{x}\in\mathcal{X}_p, \myM{u}\in\mathcal{U}_p, \boldsymbol{\lambda}\in\mathcal{K}_p.$$

where $\myM{g}(\cdot,\cdot,\cdot)$ describes the contact dynamics, and they represents terminal constraints in each walking phase. In this example we use the following \textbf{impact model}:

$$\myM{M}(\myM{v}_{next}-\myM{v}) = \myM{J}_{impulse}^T$$

$$\myM{J}_{impulse} \myM{v}_{next} = \myM{0}$$

$$\myM{J}_{c} \myM{v}_{next} = \myM{J}_{c} \myM{v}$$


\listoffigures 
\listoftables 
%%% style is changed only once, need to reopen texstudio for change %%%
%\bibliographystyle{plain} %Verzeichnis nach Autor sortiert, Referenzen numerisch
%\bibliographystyle{abbrv} %Verzeichnis nach Autor sortiert, Vornamen abgek�rzt 
%\bibliographystyle{alpha} %Verzeichnis nach Autor sortiert, Referenzen aus Autorenk�rzel
\bibliographystyle{unsrtnat} %Verzeichnis nach Zitierreihenfolge sortiert und kompatibel mit natbib package
\bibliography{tex/_bibfile}	%refer to .bib file in folder
%\nocite{*}				 % Alle Quelleneintr�ge anzeigen, auch wenn sie nicht im Text referenziert sind
\end{document}
